% Incorporación ampliativa DJ-05 del inventario REV05.
% Residencia: artículo III, después de inversión cúbica y recuperación angular.
% Propietario: output/REPARACION_ENTREGA_MULTIPLATAFORMA_20260916/ENTREGA/ES/
% PAQUETES/03_VACIO_ELECTROMAGNETICO_ES/payload/spanish_source/manuscrito/
% sections/04_respuesta_constitutiva.tex, líneas 14--65 y 139--186.
% Estatuto: FORMALIZACION_NUEVA de la diferenciación y del dominio imagen;
% la respuesta, su monotonía y su inversa cúbica son RESULTADO_RECUPERADO.

\section[Differential inversion of the constitutive response]{Differential inversion and domain of the two-sheet constitutive response}
\label{jc:sec:constitutiva}

The angular channels generated by the TPK calendar determine the positive transport
\[
T=e^{-(x+y)}P_++e^{-(x-y)}P_-,\qquad
\Omega=\{(x,y)\in\mathbb R^2:x>|y|\}.
\]
Here $P_+$ and $P_-$ are the labelled projectors of the two sheets. The incidences $90$ and $120$ act on that same transport and produce the response
\[
\mathcal V=(I-T^{90})(I-T^{120})^{-1}=r_+P_++r_-P_-.
\]
The factorization with $S=T^{30}$ determines
\begin{equation}
f(s)=\frac{1+s+s^2}{1+s+s^2+s^3},\qquad
r_\pm=f\bigl(e^{-30(x\pm y)}\bigr).
\label{jc:eq:respuesta}
\end{equation}
The factor $30$ preserves the common unit of the incidences $90=3\cdot30$ and $120=4\cdot30$. This section receives the angular coordinates previously produced; it studies regularity and inversion of their constitutive reading.

\subsection{Spectral sensitivity and logarithmic coordinates}

For $t>0$, define
\[
g(t)=\log f(e^{-30t}),\qquad \sigma(t)=g'(t),
\]
and the quotient and inverse-product coordinates
\begin{align}
B(x,y)&=\log\frac{r_-}{r_+}=g(x-y)-g(x+y),\\
V(x,y)&=\log\frac1{r_-r_+}=-g(x-y)-g(x+y).
\label{jc:eq:bv}
\end{align}
The letters $B,V$ here denote two scalar response coordinates; the sheet projectors remain in the operator $\mathcal V$.

\begin{lemma}[Regular monotonicity of the spectral reader]
The map $g$ is a strictly increasing analytic bijection from $(0,\infty)$ to $(-\log(4/3),0)$. With $s=e^{-30t}$,
\begin{equation}
\sigma(t)=
\frac{30s^3(3+2s+s^2)}{(1+s+s^2)(1+s+s^2+s^3)}>0.
\label{jc:eq:sigma}
\end{equation}
\end{lemma}
\begin{proof}
Differentiating the quotient gives
\[
f'(s)=-\frac{s^2(3+2s+s^2)}{(1+s+s^2+s^3)^2}<0,
\qquad 0<s<1.
\]
The chain rule produces $g'(t)=-30s f'(s)/f(s)$, which is \eqref{jc:eq:sigma}. The denominators are positive on the domain. The limits $f(1)=3/4$ and $f(0)=1$ give $g(0^+)=-\log(4/3)$ and $g(+\infty)=0$. Monotonicity, continuity and these limits prove bijectivity. Positivity of the derivative guarantees analytic regularity of the inverse on the open interval.
\end{proof}

\begin{theorem}[Constitutive Jacobian and local recovery]
\label{jc:thm:jacobiano} Let $a=\sigma(x-y)$ and $b=\sigma(x+y)$. The map $\Phi(x,y)=(B(x,y),V(x,y))$ satisfies
\begin{equation}
D\Phi=
\begin{pmatrix}a-b&-a-b\\-a-b&a-b\end{pmatrix},\qquad
\det D\Phi=-4ab<0.
\label{jc:eq:jacobiano}
\end{equation}
In particular, it is locally invertible throughout $\Omega$, and the differential identities $B_x=V_y$ and $B_y=V_x$ hold.
\end{theorem}
\begin{proof}
Direct differentiation of \eqref{jc:eq:bv} gives $B_x=a-b$, $B_y=-a-b$, $V_x=-a-b$ and $V_y=a-b$. The determinant is $(a-b)^2-(a+b)^2=-4ab$, strictly negative by \eqref{jc:eq:sigma}. The inverse function theorem applies at every point of the domain. The two differential identities follow from the equal entries of the matrix.
\end{proof}

\subsection{Image domain and global reconstruction}

\begin{theorem}[Global coordinates of the normalized response]
\label{jc:thm:global} Let $L=\log(4/3)$. The map $\Phi$ is an analytic diffeomorphism from $\Omega$ onto
\begin{equation}
\mathcal D=\{(B,V):0<V-B<2L,\quad0<V+B<2L\}.
\label{jc:eq:diamante}
\end{equation}
If $\psi=g^{-1}$, reconstruction is given by
\begin{align}
x&=\frac12\left[\psi\!\left(\frac{B-V}{2}\right)
                   +\psi\!\left(-\frac{B+V}{2}\right)\right],\\
y&=\frac12\left[\psi\!\left(-\frac{B+V}{2}\right)
                   -\psi\!\left(\frac{B-V}{2}\right)\right].
\label{jc:eq:inversa}
\end{align}
The chamber $x>y>0$ corresponds to the half $B<0$ of $\mathcal D$.
\end{theorem}
\begin{proof}
The change $(x,y)\mapsto(t_-,t_+)=(x-y,x+y)$ is a linear bijection between $\Omega$ and $(0,\infty)^2$. Applying $g$ to both coordinates gives an analytic bijection onto $(-L,0)^2$. If $u=g(t_-)$ and $v=g(t_+)$, then $B=u-v$ and $V=-u-v$. Its linear inverse is $u=(B-V)/2$, $v=-(B+V)/2$. The conditions $-L<u,v<0$ are exactly equivalent to \eqref{jc:eq:diamante}. Composing these three bijections gives \eqref{jc:eq:inversa} and proves the global assertion. Finally, $y>0$ is equivalent to $t_+>t_-$; monotonicity of $g$ transforms it into $v>u$, namely $B<0$.
\end{proof}

The function $\psi$ is evaluated through the cubic inverse of the reader already constructed. For $u\in(-L,0)$, take $r=e^u$ and the unique root $s\in(0,1)$ of
\[
rs^3+(r-1)(1+s+s^2)=0.
\]
Then $\psi(u)=-(\log s)/30$. Differential reconstruction and cubic reconstruction are therefore two expressions of the same constitutive inversion.

\begin{corollary}[Sheet conjugation]
The involution $(x,y)\mapsto(x,-y)$ acts on the response through $(B,V)\mapsto(-B,V)$. It preserves the product $r_+r_-$ and inverts the quotient $r_-/r_+$.
\end{corollary}
\begin{proof}
The involution exchanges $x-y$ with $x+y$ and consequently $r_-$ with $r_+$. Formula \eqref{jc:eq:bv} gives the result.
\end{proof}

\subsection{Local stability of recovery}

\begin{proposition}[Singular values and sensitivity bound]
The singular values of $D\Phi$ are $2a$ and $2b$. In particular,
\[
\bigl\|(D\Phi)^{-1}\bigr\|_2
=\frac1{2\min\{\sigma(x-y),\sigma(x+y)\}}.
\]
For $0<\delta<M<\infty$, differential recovery is uniformly bounded on $\Omega_{\delta,M}=\{(x,y):\delta\le x-y,x+y\le M\}$.
\end{proposition}
\begin{proof}
The symmetric matrix \eqref{jc:eq:jacobiano} has the orthogonal vectors $(1,-1)$ and $(1,1)$ as eigenvectors, with eigenvalues $2a$ and $-2b$. The singular values are their absolute values. The spectral norm of the inverse is the reciprocal of the smallest singular value. The positive continuous function $\sigma$ attains a positive minimum on $[\delta,M]$, providing the uniform bound.
\end{proof}

\begin{example}[Exact recovery of two rational channels]
Take $e^{-30(x-y)}=1/2$ and $e^{-30(x+y)}=1/4$. Then
\[
x=\frac{\log2}{20},\qquad y=\frac{\log2}{60},\qquad
r_-=\frac{14}{15},\qquad r_+=\frac{84}{85}.
\]
The constitutive coordinates are
\[
B=\log\frac{17}{18},\qquad V=\log\frac{425}{392}.
\]
Cubic inversion recovers $1/2$ and $1/4$, and formulas \eqref{jc:eq:inversa} restore $x,y$. At this point,
\[
a=\frac{34}{7},\qquad b=\frac{114}{119},\qquad
\det D\Phi=-\frac{15504}{833}.
\]
All rational coefficients in the example are obtained by evaluating the constitutive function and its derivative; the example verifies inversion without modifying the channel generator.
\end{example}

The transformation $(x,y)\leftrightarrow(B,V)$ describes exactly which angular information remains accessible through the two normalized readings. Dimensional realizations of permittivity, permeability, impedance and velocity are subsequently applied with their dimensional lines and declared sections. The preceding differential signature is a property of this constitutive coordinate map and retains that interpretive domain.
